\documentclass[10pt,a4paper]{amsart}
\usepackage[margin=19mm]{geometry}
\usepackage{amsmath,amssymb,mathtools}
\usepackage[scr=rsfs,cal=euler]{mathalfa}
\usepackage{xfrac}
\usepackage[unicode,hidelinks,bookmarks=false]{hyperref}
\newcommand{\ie}{i.e.\ }
\newcommand{\cf}{cf.\ }
\newcommand{\resp}{respectively\ }
\newcommand{\Subsection}{\subsection}
\providecommand{\nobreakdash}{\nobreak\hbox{-}}
\setlength{\emergencystretch}{2em}
\multlinegap0pt
\usepackage{bm}
\usepackage{bbm}
\usepackage{dsfont}
\usepackage{xspace}
\usepackage{cancel}
\usepackage{mathtools}
\usepackage{amsthm}
\makeatletter
\@ifundefined{thm}{\newtheorem{thm}{Thm}}{}
\@ifundefined{lem}{\newtheorem{lem}[thm]{Lemma}}{}
\makeatother
\def\F{\mathbb F}
\def\Z{\mathbb Z}
\def\ie{i.e.,\ }
\def\whX{\widehat{X}}
\newcommand{\Frob}{\mathrm{Frob}}
\newcommand{\Pic}{\mathrm{Pic}}
\newcommand{\ab}{\mathrm{ab}}
\title[The Pythagoras number of fields of transcendence degree $1$ ]{The Pythagoras number of fields of transcendence degree $1$ over $\mathbb{Q}$}
\author{Olivier Benoist}
\date{Source: arXiv:2506.21380v2 (CC BY 4.0)}
\begin{document}
\maketitle

\noindent\emph{Source and licence.} The Pythagoras number of fields of transcendence degree $1$ over $\mathbb{Q}$. Version arXiv:2506.21380v2 (CC BY 4.0), distributed under the Creative Commons Attribution 4.0 International licence, as recorded in the arXiv licence metadata for this exact version and shown on the abstract page https://arxiv.org/abs/2506.21380v2. Full rights notice: licenses/arxiv-2506.21380-CC-BY-4.0.txt.\par\smallskip

\noindent\textit{Editorial scope.} Counted unit: the Lem whose source label is \texttt{lemCFT} in the article (source0.tex lines 873--876, source label \texttt{lemCFT}), delivered together with the article's own complete original proof of it (source0.tex lines 878--892, one \texttt{proof} environment of 15 lines). No mathematics is written, repaired or re-derived: statement and proof are the article's own text line for line. Required context retained uncounted: lemLW (lines 820--828). Every definition, display and label of those lines is retained unchanged. Omitted from this package and not redistributed: 1--819, 829--872, 877--877, 893--1367. Nothing inside a retained excerpt is omitted or altered: every retained body is delivered exactly as the article's lines are written, so no omission receipt of any kind accompanies this package. Citations: the retained excerpts carry no citation command at all, so no bibliography entry of the article is transcribed and no reference is redistributed. Reference adaptation: none was needed and none was made; every reference inside the delivered unit resolves inside it, so no unresolved reference or citation is produced. Printed slips kept literal: no printed word, symbol, hypothesis or number of the counted statement or of its proof was altered. Layout and class: the article's own class is \texttt{amsart} with packages including amscd, amsmath, amssymb, amsfonts, bbm, calligra, xspace, fontenc, lmodern, mathtools, enumerate, enumitem, multicol, xy; the wrapper is the lane's pinned \texttt{amsart} editorial wrapper at 10pt on A4, which restates the theorem-like environments the retained text needs and adds the article's own macro definitions that the retained text needs (\texttt{\textbackslash F}, \texttt{\textbackslash Frob}, \texttt{\textbackslash Pic}, \texttt{\textbackslash Z}, \texttt{\textbackslash ab}, \texttt{\textbackslash ie}, \texttt{\textbackslash whX}), with extra wrapper packages (bm, bbm, dsfont, xspace, cancel, mathtools). The delivered numbering is the unit's own. Assets: no retained span contains a figure environment, a graphics-inclusion command, a PDF-page inclusion, a table or a TikZ picture; no image file is copied into this package, the package declares no asset dependency and no image file is redistributed with it. Licence: version arXiv:2506.21380v2 of this article is distributed under the Creative Commons Attribution 4.0 International licence, as recorded in the arXiv licence metadata for this exact version and shown on the abstract page https://arxiv.org/abs/2506.21380v2. A full rights notice is delivered at licenses/arxiv-2506.21380-CC-BY-4.0.txt. Characters: every retained excerpt is pure ASCII, so the delivered engine, XeLaTeX with its default Latin Modern text font, reports no missing character.
\begin{lem}
\label{lemLW}
Let $X$ be a geometrically integral variety of dimension $n\geq 1$ over~$\F_q$. For~$r\geq 1$, let $M_X(r)$ be the number of closed points of $X$ that are of degree~$r$ over~$\F_q$. Then there exists $K\geq 0$ such that
\begin{equation}
\label{LWclosed}
|M_X(r)-\frac{q^{nr}}{r}|\leq K\,\frac{q^{(n-\frac{1}{2})r}}{r}\textrm{\hspace{1em} for all }r\geq 1.
\end{equation}
In particular, one has $M_X(r)>0$ for all large enough integers $r\geq 1$.
\end{lem}

\begin{lem}
\label{lemCFT}
Let $X$ be a connected smooth projective curve over $\F_q$. Let~$D\subset X$ be an effective divisor. For any $d\geq 1$, the group $\Pic(X,D)^d$ is gen\-er\-at\-ed by the classes of those closed points of $X\setminus D$ whose degree over $\F_q$ is a multiple~of~$d$.
\end{lem}

\begin{proof}
After replacing $\F_q$ with its algebraic closure in $X$, we may assume that $X$ is geometrically connected. Let $G_d\subset \Pic(X,D)^d$ be the subgroup generated by 
the classes of those closed points of $X\setminus D$ whose degree over $\F_q$ is a multiple~of~$d$.
Consider the finite group $G_d^0:=G_d\cap\Pic(X,D)^0$. It follows from Lemma \ref{lemLW} that~$\deg(G_d)=d\mkern 1mu\Z$, so there is a short exact sequence $0\to G_d^0\to G_d\xrightarrow{\deg}d\mkern 1mu\Z\to 0$.

Define $H_d^0:=\rho_{(X,D)}(G_d^0)$ and let $H_d\subset \pi_1^{\ab}(X,D)$ be the closure of $\rho_{(X,D)}(G_d)$. One then has a short exact sequence $0\to H_d^0\to H_d\xrightarrow{\deg}d\mkern 1mu\widehat{\Z}\to 0$. As $H_d$ is a closed subgroup of finite index in $\pi_1^{\ab}(X,D)$, it corresponds to an abelian connected finite \'etale cover $\mu_d:\whX_d\to X$. The degree $\delta_d$ of $\mu_d$ equals the index of $H_d$ in~$\pi_1^{\ab}(X,D)$. Since~$\deg(H_d)=d\mkern 1mu\widehat{\Z}$, the algebraic closure of $\F_q$ in $\whX_d$ is isomorphic to~$\F_{q^d}$.
It follows that $\mu_d$ factors as a composition $\whX_d\to X_{\F_{q^d}}\to X$.

By Lemma \ref{lemLW}, the number of closed points of $X$ of degree $rd$ over $\F_q$ grows as~$\frac{q^{rd}}{rd}$ when $r$ goes to infinity. For any such point $x\in X$, one has $\Frob_x\in H_d$ by definition of $H_d$. This means that $x$ splits completely in $\whX_d$, \ie that the preimage of $x$ in $\whX_d$ consists of $\delta_d$ points of degree $rd$ over $\F_q$. Consequently, the number of points of $\whX_d$ of degree $rd$ over $\F_q$ grows at least as~$\delta_d\frac{q^{rd}}{rd}$ when $r$ goes to infinity. On the other hand, it follows from Lemma \ref{lemLW} (applied to the variety $\whX_d$ over $\F_{q^d}$) that this number grows as $\frac{q^{rd}}{r}$ when $r$ goes to infinity. We deduce that $\delta_d\leq d$.

The morphism $\whX_d\to X_{\F_{q^d}}$
must therefore be an isomorphism (as the degrees of the connected finite \'etale covers 
$\whX_d$ and $X_{\F_{q^d}}$ of $X$ 
are $\delta_d$ and $d$). They are thus associated with the same finite index subgroup of $\pi_1^{\ab}(X,D)$. Taking the inverse image of these subgroups by $\rho_{(X,D)}$ yields the desired equality $G_d=\Pic(X,D)^d$.
\end{proof}

\end{document}
